% Eight examples in the manner of The TeXbook: each adds one idea to the
% one before it, dangerous bends mark skippable depth, one refusal is
% shown on purpose, and the exercises' answers close the chapter.  Every
% object is named in the tensor-network community's own vocabulary
% (arXiv:2011.12127) before its spelling appears.
\section{A twenty-minute tutorial}\label{ch:tutorial}

Everything you will type here is the kernel surface: the one language the
package binds the moment it loads.  If you have met an older tenkz
document, two sentences make you current.  \tncmd{tenkzkernel} is an inert
word --- it rebinds what loading already bound --- and the spellings it
once switched on are dead, tombstoned in Appendix~\ref{ch:compatibility}.
New sources need neither.

Eight examples take you from one word of tensors to the expectation value
of a two-dimensional state, and each adds exactly one idea to the one
before it.  The objects are the ones the tensor-network literature draws
--- the matrix product state, its transfer matrix, the gauge identity,
the PEPS and its double layer --- and each is named in that language as
it arrives.  Every source is complete: paste it into any document that
loads the package and it compiles as printed.  Read each one in the same
order --- layout, picture policy, then body declarations.  This tutorial
needs only \tnenv{tenkz} and, at the equation step, \tnenv{tenkzeq}; when
a figure of your own wants a ring or another carrier this ladder does not
climb, the chooser in Section~\ref{ch:chooser} decides by topology.

\subsection{A word of tensors}

The first object is the matrix product state: one tensor $A$ per site,
with two virtual indices of bond dimension $D$ that neighbouring sites
sum over and one physical index of dimension $d$ that stays free.  Here
are three sites of it.

\begin{tnexample}[
  label={tnex:tut-word},
  formula={(A^{i_1}A^{i_2}A^{i_3})_{\alpha\beta}},
  caption={three sites of a matrix product state},
  index={The horizontal bonds are summed virtual indices.  The two stubs
    are the open matrix indices $\alpha$ and $\beta$; each upward leg is
    one physical index.},
  signature={kernel-boundary|signature=open:e, open:w, phys:n, phys:n, phys:n}]
% Formula: (A^{i_1} A^{i_2} A^{i_3})_{alpha beta}
% Ink: bonds contract; typed ports declare the open indices.
\begin{tenkz}[cols=3]
  \tn[ports={180:virtual:$\alpha$, 0:virtual, 90:physical:$i_1$}]{A} &
  \tn[ports={180:virtual, 0:virtual, 90:physical:$i_2$}]{A} &
  \tn[ports={180:virtual, 0:virtual:$\beta$, 90:physical:$i_3$}]{A}
\end{tenkz}
\end{tnexample}

You have drawn the matrix-product word $(A^{i_1}A^{i_2}A^{i_3})_{\alpha\beta}$.
\tnenv{tenkz} holds one typed
tensor network in one frame: \tnkey{cols=3} lays three cells on a wire,
each \tncmd{tn} declares an atom in the next cell, and \verb|&| advances
one site.  Adjacent cells contract by default --- those are the virtual
bonds --- so the only indices you declare are the open ones.
\tnkey{ports=} types each face and may label it; a port is the point
where an index line meets its tensor.  A typed port with no partner is
an open index --- an uncontracted one.  Open ink is declared content:
the picture exposes exactly the indices its atoms state, and the margin
quotes the boundary signature --- the list of the picture's uncontracted
indices --- that the compiler wrote to the \tnfile{.tnlog} event stream:
two virtual stubs, three physical legs, nothing you did not say.

\dbendpar Beamer reads a frame's body before typesetting it, so a chained
picture reaches the environment with its \verb|&| already an ordinary
alignment tab.  The picture reads its body as captured tokens and finds
the tab there itself, so it compiles in a plain frame; no
\texttt{[fragile]} marking is needed.

\tnexercise{ex:tut-matrix} Draw the plain matrix $M$ --- one atom, both
virtual indices open, no physical leg --- and say what boundary
signature, what list of uncontracted indices, it writes.

\subsection{Layers and side policy}

One new idea: rows are layers.  The object is the map every
matrix-product calculation applies and reapplies, the transfer channel
$\mathcal E_A(X)=\sum_i A^iXA^{i\dagger}$ --- a ket layer over a bra
layer, with the matrix $X$ closing the input side.

\begin{tnexample}[
  formula={\mathcal E_A(X)=\sum_i A^i X A^{i\dagger}},
  caption={the channel applied to $X$},
  index={$X$ closes the east input pair.  The west pair remains open because
    the output is a matrix; the vertical connection sums the physical
    index $i$.},
  signature={kernel-boundary|signature=open:w, open:w}]
% Formula: E_A(X) = sum_i A^i X A^{i dagger}.
% Boundary: the input closes east; the matrix output stays
% open west.
\begin{tenkz}[rows={ket,bra}, cols=1, west=open, east={cup=$X$}]
  \tn[label pos=90]{A} \\
  \tn[label pos=270]{$\overline{A}$}
\end{tenkz}
\end{tnexample}

\tnkey{rows=} declares a ket layer over a bra layer, and the frame
contracts their facing physical ports --- the vertical bond that sums $i$.
The conjugate is label mathematics: $\overline{A}$ is a name, not a flag.
The sides then state policy from a four-word alphabet.  \tnval{open} mints
legs and puts them in the signature; \tnval{none}, the default, is a side
with no indices at all; \tnval{trace} returns a row to itself; \tnval{cup}
closes adjacent layers onto each other.  Here \tnkey{east=}\tnval{\{cup=\$X\$\}}
closes the pair through the matrix $X$, and the west pair stays open, so
the picture is a map on virtual matrices --- which is exactly what its
two-entry signature says.

\dbendpar The labelled cup is sugar: it expands to the side word
\tnval{cup} plus a ring atom standing halfway along the generated cup
wire, so the kernel record stream never contains a ``labelled cup'',
only a wire named \tnval{cup-1-2} and an atom on it.  Every sugar row in
the registry states its expansion this way, and \tncmd{tnset}\tnval{\{strict\}}
rejects sugar outright.

\tnexercise{ex:tut-transfer} Draw the transfer matrix
$E=\sum_i A^i\otimes\overline{A^i}$ of the state --- the map on the
$D^2$-dimensional doubled virtual space: ket over bra, the physical pair
contracted, all four virtual ends open.

\tnexercise{ex:tut-closed} A contraction loop with no tensor on it ---
the spelling \verb|\tnwire[closed]{(1,1)}{(2,1)}| --- refuses to
compile.  Why?

\subsection{The boundary alphabet, three ways}

One row of matrices is three objects of the literature: a word with
trivial boundary, the open-boundary matrix product
$(M_1M_2M_3)_{\alpha\beta}$, and the periodic trace
$\operatorname{tr}(M_1M_2M_3)$ --- the coefficient of a matrix product
state with open and with periodic boundary conditions.  The pair
spelling \tnkey{boundary=} states \tnkey{west=} and \tnkey{east=} at
once, and its three words draw all three objects from one row:

\begin{tnmultiples}[
  formula={(M_1M_2M_3)_{\alpha\beta}},
  caption={one row, three boundary conditions},
  index={Left to right: the outer bonds close trivially and nothing is
    exposed; both ends open and the row is a matrix; the trace returns the
    row to itself and the picture is a scalar.  The quoted signature is the
    open panel's; the other two are empty.},
  signature={kernel-boundary|signature=open:e, open:w},
  variants={{boundary=none}{boundary=none},
            {boundary=open}{boundary=open},
            {boundary=periodic}{boundary=periodic}}]
% Formula: M_1 M_2 M_3 under three side policies.
\begin{tenkz}[rows={wire}, cols=3, variant]
  \tn{$M_1$} & \tn{$M_2$} & \tn{$M_3$}
\end{tenkz}
\end{tnmultiples}

One family, one figure: the three panels share every token except the
word under them.  \tnval{none} is the explicit spelling of the default
--- a side that says nothing exposes nothing, the way a bond that closes
trivially at the chain's outer ends deserves no ink.  \tnval{open} turns
the row into the matrix $M_1M_2M_3$, and \tnval{periodic} (sugar for
\tnkey{west=}\tnval{trace}, \tnkey{east=}\tnval{trace}) closes the row on
itself.  Same atoms, same bonds; only the boundary signature moves.

\dbendpar The default is a statement, not an omission.  The deleted
0.7 dialect drew west and east stubs unbidden and sealed them with
\tnval{none}; the kernel mints a leg only on declaration, so
\tnval{none} now changes nothing and merely says so.  That reading died
with the old front end --- if you sealed stubs once, stop: there are no
stubs to seal.

\tnexercise{ex:tut-trace} Draw the coefficient of the periodic matrix
product state --- the scalar $\operatorname{tr}(M_1M_2M_3)$ --- and say
what its boundary signature is.

\subsection{A mark over a run}

Blocking is the coarse-graining step of every renormalization argument
in the matrix-product literature: three sites read as one blocked tensor
$A^{[3]}$.  The contraction does not change; a mark speaks over it.

\begin{tnexample}[
  formula={A^{[3]}=A^{i_1}A^{i_2}A^{i_3}},
  caption={blocking three sites},
  index={The bracket speaks from the south, past the site names and their
    ink.  The signature is Example \ref{tnex:tut-word}'s, unchanged: a
    mark owns no topology.},
  signature={kernel-boundary|signature=open:e, open:w, phys:n, phys:n, phys:n}]
% Formula: the blocked word A^{[3]}; the mark adds no index.
\begin{tenkz}[cols=3]
  \tn[ports={180:virtual:$\alpha$, 0:virtual, 90:physical:$i_1$}]{A} &
  \tn[ports={180:virtual, 0:virtual, 90:physical:$i_2$}]{A} &
  \tn[ports={180:virtual, 0:virtual:$\beta$, 90:physical:$i_3$}]{A}
  \tnmark[form=bracket]{(1,1) .. (1,3)}{$A^{[3]}$}
\end{tenkz}
\end{tnexample}

\tncmd{tnmark} annotates a target without touching the contraction.  The
target here is a selector: \tnval{(1,1)} names a cell of the frame, and
the two-dot range \tnval{(1,1) .. (1,3)} names the records between two
cells in the frame's own order.  A bracket that names no side speaks from
the south, where a brace under a span belongs.  Compare the margin with
Example~\ref{tnex:tut-word}: the signature has not moved, because marks
never own topology --- delete every mark and no contraction changes.

\dbendpar The range is written with two dots, never a hyphen, because
generated record names carry hyphens of their own: a cluster member is
named \tnval{a-2-2}, and a hyphen range could not be told from it.

\tnexercise{ex:tut-enclosure} Restate the same blocked run as a tinted
enclosure whose label sits below the contour.

\subsection{An equation under audit}

The fifth idea is the one the others were preparing: whole pictures
compose into an equation, and the equation is an audit scope.  The
identity is the gauge freedom of the matrix product state --- two
tensors, $B^i=XA^iX^{-1}$, generating the same state.

\begin{tnexample}[
  label={tnex:tut-eq},
  formula={B^i=XA^iX^{-1}},
  caption={a gauge identity, audited},
  index={The rings are matrices on the wire and add no indices.  Both
    panels expose one west stub, one east stub, and one physical leg, so
    the audit records the relation as equal.},
  signature={check|scope=1|relation=1|result=equal|signature=open:e, open:w, phys:n}]
% Formula: B^i = X A^i X^{-1}; one relation, audited.
\[
\begin{tenkzeq}[check=signature]
  \begin{tenkz}[cols=1]
    \tn[ports={180:virtual, 0:virtual, 90:physical:$i$}]{B}
  \end{tenkz}
  =
  \begin{tenkz}[cols=3]
    \tn[skin=ring, ports={180:virtual, 0:virtual}]{X} &
    \tn[ports={180:virtual, 0:virtual, 90:physical:$i$}]{A} &
    \tn[skin=ring, ports={180:virtual, 0:virtual}]{X^{-1}}
  \end{tenkz}
\end{tenkzeq}
\]
\end{tnexample}

\tnenv{tenkzeq} places its panels around the mathematics between them ---
the equals sign is ordinary TeX --- under one shared metric, so a glyph
is one size on both sides of an equality.  Then it audits: every panel
carries a boundary signature, and a relation requires its two sides to
expose the same one.  \tnkey{check=}\tnval{signature} names that audit
(it is also the default: writing it states what is already running), and
the margin quotes the \tnval{check} event the equation wrote beside each
panel's boundary line.  A diagrammatic identity in this language is not a
juxtaposition that looks right; it is a comparison that ran.

What does a failed comparison look like?  Drop the physical leg of $B$
and the compiler answers:

\begin{tnrefusal}[
  caption={the audit catches a dropped leg},
  index={The left panel now exposes two indices, the right three.  The
    correction is one token: restore the port 90:physical:$i$ on $B$.},
  diagnostic={[TKZ-EQ-SIGNATURE] sides 1 and 2 expose unequal boundaries:
    (open:e, open:w) versus (open:e, open:w, phys:n).}]
% Refused: B lost its physical leg, so the relation's two
% sides no longer expose the same boundary.
\[
\begin{tenkzeq}[check=signature]
  \begin{tenkz}[cols=1]
    \tn[ports={180:virtual, 0:virtual}]{B}
  \end{tenkz}
  =
  \begin{tenkz}[cols=3]
    \tn[skin=ring, ports={180:virtual, 0:virtual}]{X} &
    \tn[ports={180:virtual, 0:virtual, 90:physical:$i$}]{A} &
    \tn[skin=ring, ports={180:virtual, 0:virtual}]{X^{-1}}
  \end{tenkz}
\end{tenkzeq}
\]
\end{tnrefusal}

Read refusals as language, not as accidents: the code names the check
that failed, and the message prints both signatures so the missing entry
is visible at once.  The correction is the port list of
Example~\ref{tnex:tut-eq}'s~$B$.

\ddbendpar Relations are one of three joiner classes.  A sum also
requires equal signatures; a product instead contracts the facing cut,
east entries of the left factor cancelling west entries of the right,
and exposes the rest.  A relation whose two sides are known, unequal, and
equal anyway for a reason the diagram does not draw carries a recorded opt-out,
\tnval{off=\{<relation>: <reason>\}}, written to the event stream beside
the waiver.  The audit cannot be switched off silently.

\tnexercise{ex:tut-checkword} Why does
\tnkey{check=}\tnval{\{signature=false\}} refuse?

\subsection{A sheet of tensors}

The chain taught the whole alphabet; the second dimension only adds a
frame.  The object is the projected entangled pair state (PEPS), the
wavefunction of a two-dimensional system: one five-index tensor $A$ per
lattice site, four virtual indices of bond dimension $D$ contracted with
the four neighbours, one physical index of dimension $d$ standing free
above the sheet, and the tensor $A$ generating the state
$\lvert\psi(A)\rangle$.

\begin{tnexample}[
  label={tnex:tut-peps},
  formula={\lvert\psi(A)\rangle},
  caption={a PEPS on a $3\times3$ patch},
  index={Every dot is the same tensor $A$.  The in-sheet bonds are the
    contracted virtual indices; each transverse leg is one physical
    index, and the signature lists the nine of them.},
  signature={kernel-boundary|signature=phys:n, phys:n, phys:n, phys:n,
    phys:n, phys:n, phys:n, phys:n, phys:n}]
% Formula: the PEPS |psi(A)> generated by one tensor A.
% Ink: in-sheet bonds contract; each transverse leg is one
% physical index.
\begin{tenkz}[lattice={3x3}, frame=plane, physical=up]
  \tn{} & \tn{} & \tn{} \\
  \tn{} & \tn{} & \tn{} \\
  \tn[label pos=225]{A} & \tn{} & \tn{}
\end{tenkz}
\end{tnexample}

Two spellings put the sheet on the page.  \tnkey{lattice=}\tnval{\{3x3\}}
is sugar for three wire rows of three columns, and the cells contract
with their neighbours exactly as chain cells do --- those are the virtual
bonds of the sheet.  \tnkey{frame=}\tnval{plane} projects the grid, and
\tnkey{physical=}\tnval{up} grows one physical leg per site along the
plane's transverse axis --- the direction that stands vertically on the
page and belongs to the frame itself.  No side policy is stated, so no
virtual index is exposed: the patch has open boundary conditions, and its
boundary signature is the nine physical legs of the wavefunction.

\dbendpar A plane owns three directions, and only two lie in the sheet.
\tnval{ports=\{90:physical\}} on a plane atom is the in-sheet north face,
not the transverse leg; the physical axis is reached only through the
picture policy.  A PEPS tensor is four numeric ports plus the policy ---
never a fifth angle compensating for the projection.

\tnexercise{ex:tut-pepstensor} Draw one PEPS tensor by itself --- the
five-index $A^{i}_{\,uldr}$: four open virtual ports labelled $u$, $l$,
$d$, $r$ in the sheet, and the physical leg above.

\subsection{One sheet, four readings}

The same sheet is four objects of the two-dimensional literature, and
one enum word turns between them: a fully contracted network is a
partition function $Z$; legs up is the state $\lvert\psi\rangle$; legs
down is its conjugate $\langle\psi\rvert$; legs both ways is a projected
entangled pair operator $O$, the PEPO.

\begin{tnmultiples}[
  formula={Z,\;\lvert\psi\rangle,\;\langle\psi\rvert,\;O},
  caption={one sheet, four readings},
  index={Left to right: no physical legs, a fully contracted scalar;
    legs up, the state; legs down, its conjugate; legs both ways, the
    PEPS operator.  The quoted signature is the legs-up panel's.},
  signature={kernel-boundary|signature=phys:n, phys:n, phys:n, phys:n},
  variants={{physical=none}{physical=none},
            {physical=up}{physical=up},
            {physical=down}{physical=down},
            {physical=updown}{physical=updown}}]
% Formula: Z, |psi>, <psi|, O -- one sheet, four leg policies.
\begin{tenkz}[lattice={2x2}, frame=plane, variant]
  \tn{} & \tn{} \\
  \tn{} & \tn{}
\end{tenkz}
\end{tnmultiples}

One family, one figure again: the four panels share every token except
the word under them, and the family is enumerated in full.  The
signature tells a ket from a bra before any render is looked at: a leg
above the sheet writes \tnval{phys:n}, a leg below writes \tnval{phys:s},
so the record stream distinguishes the state from its dual, and the PEPO
panel exposes both lists at once.

\tnexercise{ex:tut-torus} A PEPS with periodic boundary conditions lives
on a torus.  The one word \tnkey{surface=}\tnval{torus} closes all four
sides of the patch --- sugar for \tnval{trace} on each of \tnkey{west=},
\tnkey{east=}, \tnkey{north=}, and \tnkey{south=}.  What does it change
in the boundary signature of Example~\ref{tnex:tut-peps}, and what in
the network?

\subsection{The double layer}

Every expectation value in the PEPS literature is computed in the double
layer: the conjugate sheet laid over the state, each site's physical
index contracted with its conjugate copy.  The norm
$\langle\psi|\psi\rangle$ is the double layer with nothing inserted.

\begin{tnexample}[
  label={tnex:tut-doublelayer},
  formula={\langle\psi|\psi\rangle},
  caption={the norm as a double layer},
  index={Two $2\times2$ sheets in one frame: $A$ marks the ket sheet,
    $\overline{A}$ the bra copy above it.  Each short diagonal is one
    ket--bra physical contraction; nothing stays open, and the empty
    signature says scalar.},
  signature={kernel-boundary|signature=}]
% Formula: <psi|psi> -- the bra sheet contracted onto the
% ket sheet through every physical index.
\begin{tenkz}[lattice={2x2}, planes, bonds=none]
  \tn[at=(2,1,1), label pos=225]{A}
  \tn[at=(1,2,2), label pos=45]{$\overline{A}$}
  % ket-sheet virtual bonds
  \tnwire{(1,1,1)}{(1,2,1)} \tnwire{(2,1,1)}{(2,2,1)}
  \tnwire{(1,1,1)}{(2,1,1)} \tnwire{(1,2,1)}{(2,2,1)}
  % bra-sheet virtual bonds
  \tnwire{(1,1,2)}{(1,2,2)} \tnwire{(2,1,2)}{(2,2,2)}
  \tnwire{(1,1,2)}{(2,1,2)} \tnwire{(1,2,2)}{(2,2,2)}
  % ket-bra physical pairings
  \tnwire{(1,1,1)}{(1,1,2)} \tnwire{(1,2,1)}{(1,2,2)}
  \tnwire{(2,1,1)}{(2,1,2)} \tnwire{(2,2,1)}{(2,2,2)}
\end{tenkz}
\end{tnexample}

\tnval{planes} is sugar for one plane frame with a two-member basis ---
a ket sheet and a bra sheet --- so both copies live in one frame and
every site has a member address: \tnval{(r,c,1)} is the ket copy of cell
\tnval{(r,c)}, \tnval{(r,c,2)} the bra copy.  A multi-member basis draws
no bond it is not told --- \tnkey{bonds=}\tnval{none} says so --- and the
body declares each one: eight virtual bonds, sheet by sheet, then four
ket--bra pairings.  A pairing carries the type neither of its ends can
state, because a member address names a lattice place and a place is not
a port; the frame knows its transverse axis, types the wire physical,
and the record stream says so.

\dbendpar The two sheets interleave in projection, and several bonds
pass close by sites of the other sheet.  Nothing contracts there:
coordinate coincidence neither identifies records nor creates an edge,
and the only contractions are the twelve wires the body declares.

\tnexercise{ex:tut-expectation} Insert a one-site operator.  The
expectation value $\langle\psi\rvert O\lvert\psi\rangle$ stands the
ring $O$ on one ket--bra pairing: name the pairing wire and put the
operator halfway along it.

\subsection{Answers to the exercises}

\begin{tnexample}[
  answer={ex:tut-matrix},
  formula={M_{\alpha\beta}},
  index={Two typed ports with no partners: the matrix keeps two open
    virtual indices and nothing else.},
  signature={kernel-boundary|signature=open:e, open:w}]
% Formula: the matrix M with both virtual indices open.
\begin{tenkz}[cols=1]
  \tn[ports={180:virtual:$\alpha$, 0:virtual:$\beta$}]{M}
\end{tenkz}
\end{tnexample}

\begin{tnexample}[
  answer={ex:tut-transfer},
  formula={E=\sum_i A^i\otimes\overline{A^i}},
  index={The frame contracts the ket-bra physical pair; all four virtual
    ends stay open, so $E$ is a map on the doubled space.},
  signature={kernel-boundary|signature=open:e, open:e, open:w, open:w}]
% Formula: E = sum_i A^i tensor conj(A^i).
% Labels: automatic placement -- every face of the ket atom except the
% north one carries ink, so its name stands there.
\begin{tenkz}[rows={ket,bra}, cols=1, west=open, east=open]
  \tn{A} \\
  \tn{$\overline{A}$}
\end{tenkz}
\end{tnexample}

\tnanswer{ex:tut-closed} A closed wire is a smooth cycle, so it has no
ends to give; every other wire has exactly two.  The spelling states
both at once and is refused with the coded error
\tnval{[TKZ-LANG-WIRE-ARITY]}: \emph{a closed wire takes no positional
ends; every other wire takes exactly two}.  Keep \tnval{closed} for a
free loop, or keep the two ends for a bond.

\begin{tnexample}[
  answer={ex:tut-trace},
  formula={\operatorname{tr}(M_1M_2M_3)},
  index={The trace returns the row to itself; no index is left open, and
    the empty signature says scalar.},
  signature={kernel-boundary|signature=}]
% Formula: tr(M_1 M_2 M_3); sugar: boundary=periodic
% expands to west=trace, east=trace.
\begin{tenkz}[rows={wire}, cols=3, boundary=periodic]
  \tn{$M_1$} & \tn{$M_2$} & \tn{$M_3$}
\end{tenkz}
\end{tnexample}

\begin{tnexample}[
  answer={ex:tut-enclosure},
  formula={A^{[3]}},
  index={The enclosure strokes the offset hull of the selection and tint
    lays ink over the paper inside it; the physical legs pierce the
    contour, and the signature still does not move.},
  signature={kernel-boundary|signature=open:e, open:w, phys:n, phys:n, phys:n}]
% Formula: the blocked word inside one tinted contour.
\begin{tenkz}[cols=3]
  \tn[ports={180:virtual:$\alpha$, 0:virtual, 90:physical:$i_1$}]{A} &
  \tn[ports={180:virtual, 0:virtual, 90:physical:$i_2$}]{A} &
  \tn[ports={180:virtual, 0:virtual:$\beta$, 90:physical:$i_3$}]{A}
  \tnmark[form=enclosure, tint,
          label pos=270]{(1,1) .. (1,3)}{$A^{[3]}$}
\end{tenkz}
\end{tnexample}

\tnanswer{ex:tut-checkword} The audit refuses words it does not know
rather than passing over them, because a spelling the audit ignores is a
comparison you believe is running and is not.  The coded error is
\tnval{[TKZ-EQ-CHECK-WORD]}: \emph{the audit does not know
`signature=false'; check= states the audit (signature) and its recorded opt-outs
(off=\{<relation>: <reason>\})}.  The word \tnval{signature} names the
audit; it is not a switch and therefore takes no value.

\begin{tnexample}[
  answer={ex:tut-pepstensor},
  formula={A^{i}_{\,uldr}},
  index={Four open in-sheet virtual indices and one transverse physical
    leg: the five-index PEPS tensor.  The two receding faces expose
    their in-sheet bearings in the signature.},
  signature={kernel-boundary|signature=open:233.130102,
    open:53.130102, open:e, open:w, phys:n}]
% Formula: the five-index PEPS tensor A^i_{uldr}.
\begin{tenkz}[rows={wire}, cols=1, frame=plane,
              bonds=none, physical=up]
  \tn[ports={0:virtual:$r$, 90:virtual:$u$,
             180:virtual:$l$, 270:virtual:$d$},
      label pos=315]{A}
\end{tenkz}
\end{tnexample}

\tnanswer{ex:tut-torus} Nothing moves in the signature, and six wires
appear in the network.  Each \tnval{trace} closes virtual bonds the
patch never exposed, so the boundary signature still reads nine physical
legs; the network gains one wrap wire per row and per column ---
\tnval{wrap-1} through \tnval{wrap-3} and \tnval{wrap-col-1} through
\tnval{wrap-col-3} --- closing the sheet into a torus.  The signature is
the boundary, not the topology: the open patch and the periodic state
expose one list and are different networks.

\begin{tnexample}[
  answer={ex:tut-expectation},
  formula={\langle\psi\rvert O\lvert\psi\rangle},
  index={The ring stands halfway along the named ket--bra pairing: a
    one-site operator between the state and its conjugate.  The
    signature stays empty --- the expectation value is a scalar.},
  signature={kernel-boundary|signature=}]
% Formula: <psi|O|psi> -- O on one site's physical index.
\begin{tenkz}[lattice={2x2}, planes, bonds=none]
  \tnwire{(1,1,1)}{(1,2,1)} \tnwire{(2,1,1)}{(2,2,1)}
  \tnwire{(1,1,1)}{(2,1,1)} \tnwire{(1,2,1)}{(2,2,1)}
  \tnwire{(1,1,2)}{(1,2,2)} \tnwire{(2,1,2)}{(2,2,2)}
  \tnwire{(1,1,2)}{(2,1,2)} \tnwire{(1,2,2)}{(2,2,2)}
  \tnwire{(1,1,1)}{(1,1,2)} \tnwire{(1,2,1)}{(1,2,2)}
  \tnwire{(2,1,1)}{(2,1,2)}
  \tnwire[name=po]{(2,2,1)}{(2,2,2)}
  \tn[skin=ring, at=on po 0.5]{O}
\end{tenkz}
\end{tnexample}
